\chapter{Storage: Volumes, Claims, and State on a Single Node}
\label{ch:storage}

\section{The abstraction stack}

Containers' filesystems die with them. Kubernetes separates the
\emph{request} for storage from its \emph{provision} through three
objects:

\begin{description}
  \item[PersistentVolume (PV)] an actual piece of storage --- a
  directory, an EBS volume, an NFS export.
  \item[PersistentVolumeClaim (PVC)] a workload's request: ``5\,Gi,
  ReadWriteOnce.'' The claim is what pods mount; it binds to a PV.
  \item[StorageClass] the factory: how to make a PV on demand when a
  claim appears (\emph{dynamic provisioning}).
\end{description}

The indirection is portability: \code{postgres-course.yaml} says only
``I need 5\,Gi''; \emph{how} that becomes disk is the cluster's
business. The same manifest gets a local directory on k3s and an EBS
volume on EKS.

\section{local-path on k3s}

k3s bundles the \code{local-path} StorageClass: a PVC becomes a
directory under \code{/var/lib/rancher/k3s/storage/} on the node's SSD.
Simple, fast, and honest about its limits --- the data is on \emph{that
node's disk}. On multi-node clusters that pins the pod to the node; on
this single-node cluster it just means what it says.

\begin{hardware}
local-path provides \emph{persistence} (data survives pod and node
restarts), not \emph{durability} (one SSD failure loses everything) and
not \emph{backup} (a bad \code{DELETE} replicates to nowhere).
The honest missing piece is a CronJob running \code{pg\_dump} to MinIO
or off-box. Rule of thumb worth keeping: a database without tested
restores is a cache.
\end{hardware}

\section{Two claim styles in the manifests}

MinIO (a Deployment with one replica) uses a standalone PVC; the
databases use \code{volumeClaimTemplates} inside their StatefulSets ---
a PVC \emph{minted per replica} and, crucially, \emph{re-bound to the
replacement pod} rather than a second live pod. That distinction is
the storage half of Chapter~15's corruption pitfall.

\begin{lstlisting}[language=yaml]
volumeClaimTemplates:
  - metadata: { name: data }
    spec:
      accessModes: [ReadWriteOnce]     (*@\co{1}@*)
      resources: { requests: { storage: 5Gi } }
# in the pod template:
volumeMounts:
  - { name: data, mountPath: /var/lib/postgresql/data }
env:
  - { name: PGDATA, value: /var/lib/postgresql/data/pgdata }  (*@\co{2}@*)
\end{lstlisting}

\begin{description}
  \item[\co{1}] \code{ReadWriteOnce} = mountable read-write by one
  \emph{node} at a time --- fits a database exactly. (\code{RWX} ---
  many nodes --- needs NFS-like storage and is the mode people want for
  shared uploads before discovering object storage, Chapter~9, is the
  better answer.)
  \item[\co{2}] The PGDATA trick, a genuine classic: the Postgres image
  refuses to initialize into a non-empty directory, and ext4 volumes
  arrive with \code{lost+found}. Pointing PGDATA one directory deeper
  sidesteps it. You will meet this exact pattern in
  StackOverflow-years' worth of Postgres-on-Kubernetes setups.
\end{description}

\begin{pitfall}[deleted the StatefulSet, data still there --- or gone]
Two surprises in one. First: deleting a StatefulSet does \emph{not}
delete its PVCs --- by design, data outlives orchestration mistakes;
re-create the StatefulSet and it re-binds. Second, the reverse trap:
deleting the PVC while experimenting \emph{does} delete the PV and its
directory (reclaimPolicy \code{Delete}, the dynamic-provisioning
default). The asymmetry is worth memorizing before you clean up a
namespace: workload objects are disposable; claims are where the data
is.
\end{pitfall}

\begin{handson}[see through the abstraction]
\begin{lstlisting}[language=bash]
kubectl -n ts get pvc                       # claims and their PVs
kubectl get pv                              # cluster-scoped, note the paths
ssh ts-server 'sudo ls /var/lib/rancher/k3s/storage/'
# pvc-...  one directory per PV -- your databases, as plain files
\end{lstlisting}
Demystification is the point: a PVC on this cluster is a directory.
On EKS the same two kubectl commands would show EBS volume IDs ---
same objects, different factory.
\end{handson}

\section{Outside the project}

On managed Kubernetes, StorageClasses map to cloud disks (EBS/gp3 on
EKS, Persistent Disks on GKE) with snapshots, encryption and resizing
--- \code{VolumeSnapshot} objects make backup part of the same
declarative model. The industrial pattern this chapter points toward,
though, is subtraction: teams increasingly keep \emph{databases outside
the cluster} (RDS, Cloud SQL) and let Kubernetes run only stateless
things --- trading Chapter~15's StatefulSet care for a managed bill.
Both sides of that trade are now visible to you.

\begin{quiz}
\begin{enumerate}
  \item PV, PVC, StorageClass: which does a pod mount, which is
  cluster-scoped, and what does the indirection buy this project's
  manifests on a future EKS?
  \item Why is \code{ReadWriteOnce} sufficient for Postgres, and what
  kind of workload pushes people toward RWX --- usually wrongly?
  \item Explain both directions of the deletion asymmetry: StatefulSet
  gone/PVC stays, PVC gone/data gone.
  \item What do persistence, durability, and backup each mean, and
  which does local-path provide?
\end{enumerate}
\end{quiz}
